\section{Experimental Setup}
\label{sec:setup}

\subsection{Hardware and Software}
All experiments were run on a single workstation with an NVIDIA GeForce RTX
4060 Laptop GPU (8\,GB), CUDA 12.6, PyTorch 2.13 and Python 3.11, using the
Hugging Face Transformers library \cite{wolf2020transformers}.  The
\texttt{mdeberta-v3-base} encoder has $278.2$M parameters.  The standard Hugging
Face head (a 768-d pooler plus a 3-way linear layer) adds $0.59$M, so the
standard-head systems (M1, M6) total $278.8$M; the multi-view head adds
$1.19$M, so the custom-head systems (M2--M5) total $279.4$M.  The largest gap
between any two systems is thus $+0.21\%$ of the parameters, and we treat them as
of essentially identical size.

\subsection{Training Protocol}
All systems share one configuration: AdamW \cite{loshchilov2019adamw} with weight
decay 0.01, a linear schedule with warmup $0.05$--$0.1$, batch size 16 with
gradient accumulation~2 (effective 32), gradient clipping at norm 1.0, sequence
length 256 tokens, bf16 training (fp32 for the latency benchmark), and up to
15--30 epochs depending on the recipe.  Early stopping monitors validation
macro-F1 with patience 6 and restores the best checkpoint.  The systems differ
only in the encoder initialisation, the number of stages, the loss sequence and
the class-weighting scheme (Sect.~\ref{sec:chain}); the test split is evaluated
exactly once per system with the checkpoint selected on validation.
Table~\ref{tab:training-details} lists the resolved configuration and the frozen
checkpoint (best epoch and its validation macro-F1) for each system.  Every system
reads the same three split files (verified by SHA-256 digest) and label mapping,
and all class weights and thresholds are estimated on training and validation only,
never on test; only the head, learning rates, number of stages and loss vary.

% AUTO-GENERATED by scripts/06_analyze_significance.py - do not edit by hand.
\begin{table}[t]
\centering
\small
\setlength{\tabcolsep}{4pt}
\caption{Training configuration and the validation checkpoint that was frozen for the test pass (M1--M6, Sect.~\ref{sec:chain}). AdamW, linear schedule, effective batch size 32 and sequence length 256 are shared by all systems. ``init`` is the encoder the first stage starts from; ``best ep.`` is the epoch whose validation macro-F1 was highest (it can lie in an earlier stage).}
\label{tab:training-details}
\begin{tabular}{lcccccc}
\toprule
System & init & stages & loss & class w. & best ep. & val M-F1 \\
\midrule
M1 & pre-trained & 1 & CE & none & 9 & 0.6961 \\
M2 & pre-trained & 1 & CE & none & 14 & 0.6957 \\
M3 & chained & 2 & CE $\to$ focal & inv-sqrt & 15\,(s1) & 0.7089 \\
M4 & M3 best & 3 & concept $\to$ CE & manual & 1\,(s1) & 0.6942 \\
M5 & M3 stage-1 & 1 & CE & inv-sqrt & 3 & 0.7072 \\
M6 & M1 best & 1 & focal & inv-sqrt & 4 & 0.6928 \\
\bottomrule
\end{tabular}
\end{table}


\subsection{Metrics}
Because of the imbalance described in Sect.~\ref{sec:imbalance}, we treat
\emph{macro-averaged F1} as primary and report accuracy only as a diagnostic,
alongside per-class precision, recall and F1 and macro ROC-AUC / average precision
(AP), since precision-recall summaries are more trustworthy than ROC under
imbalance \cite{saito2015precision}.  Calibration uses the top-label expected
calibration error (ECE, 10 bins); significance uses a percentile bootstrap
(macro-F1) and the exact McNemar test (accuracy).
